% =====================================================================
%  4.4 Drøfting av funksjoner — fasit
%  \fk{overskrift}   \fo{oppgavenummer} svar
% =====================================================================

\fk{Del 1 — Fortegnslinje for $f'$ og ekstremalpunkter}
\fo{1.1} a) $f'(x) = 6x^2 - 6$ \ b) $f'(x) = 6(x - 1)(x + 1)$ \ c) $x = -1$ og $x = 1$
\fo{1.2} a) Voksende på $\langle\leftarrow, -2]$ og $[1, \rightarrow\rangle$. Avtagende på $[-2, 1]$.
b) Toppunkt i $x = -2$ ($f'$ går fra $+$ til $-$). Bunnpunkt i $x = 1$ ($f'$ går fra $-$ til $+$).
\fo{1.3} a) $f'(x) = 2x - 6$, $x = 3$ \ b) $f'$ skifter fra $-$ til $+$: bunnpunkt \ c) $f(3) = -4$, punktet er $(3, -4)$
\fo{1.4} a) $f'(x) = 3x^2 + 6x - 9 = 3(x + 3)(x - 1)$. Fortegn: $+$ for $x < -3$, $-$ for $-3 < x < 1$, $+$ for $x > 1$.
b) Voksende på $\langle\leftarrow, -3]$ og $[1, \rightarrow\rangle$. Avtagende på $[-3, 1]$.
c) Toppunkt $(-3, 27)$ og bunnpunkt $(1, -5)$.
\fo{1.5} a) $g(1) = 1 - 3 - 9 + 11 = 0$. Divisjonen (skriv produktet med motsatt fortegn og legg sammen, det er det samme som å trekke fra):
{\small\polylongdiv{x^3-3x^2-9x+11}{x-1}}
gir $g(x) = (x - 1)(x^2 - 2x - 11)$. Andregradsformelen gir $x = 1 \pm 2\sqrt{3}$.
Nullpunkter: $x = 1$, $x = 1 - 2\sqrt{3} \approx -2,46$ og $x = 1 + 2\sqrt{3} \approx 4,46$.
b) $g'(x) = 3x^2 - 6x - 9 = 3(x + 1)(x - 3)$, fortegn $+, -, +$. Toppunkt $(-1, 16)$ og bunnpunkt $(3, -16)$.
c) Skissen går gjennom $(-2,46;\,0)$, $(-1, 16)$, $(0, 11)$, $(1, 0)$, $(3, -16)$ og $(4,46;\,0)$.
\fo{1.6} $f'(x) = 4x^3 - 12x^2 = 4x^2(x - 3)$, så $x = 0$ eller $x = 3$. Fortegnslinje: $x^2 \ge 0$ skifter ikke fortegn i $0$, og $x - 3$ er negativ for $x < 3$ og positiv for $x > 3$.
Så $f'$ er negativ på begge sider av $x = 0$: terrassepunkt $(0, 0)$. I $x = 3$ går $f'$ fra $-$ til $+$: bunnpunkt $(3, -27)$.
\fo{1.7} Feilen er overgangen til «$f'$ er negativ for $x < -2$». Fortegnet er $+, -, +$: for eksempel er $f'(-3) = 3(-7)(-1) = 21 > 0$.
Riktig: $f'$ er positiv for $x < -2$, negativ for $-2 < x < 4$ og positiv for $x > 4$. Toppunkt $(-2, 28)$ og bunnpunkt $(4, -80)$.
\fo{1.8} a) Noen ganger. For $f(x) = 2x$ er $f'(x) = 2 > 0$. For $f(x) = x^3$ er $f$ voksende, men $f'(0) = 0$.
b) Alltid. Med $f(x) = px^2 + qx + r$ ($p \ne 0$) er $f'(x) = 2px + q$, som er null for nøyaktig én $x$-verdi, og en førstegradsfunksjon skifter fortegn der. Det gir ett topp- eller bunnpunkt.
c) Noen ganger. $f(x) = x^2$ og $g(x) = x^2 + 5$ har samme $f'$ og samme fortegnslinje, men er ikke like. Hvis $f = g$, er fortegnslinjene like.
\fo{1.9} a) Ett svar: $f'(x) = 6(x + 1)(x - 2) = 6x^2 - 6x - 12$. Da gir $6x^2$ leddet $2x^3$, $-6x$ gir $-3x^2$ og $-12$ gir $-12x$, så $f(x) = 2x^3 - 3x^2 - 12x$. Faktoren $6$ er valgt for å unngå brøker. Fortegn for $f'$: $+, -, +$. Toppunkt $(-1, 7)$ og bunnpunkt $(2, -20)$. Flere svar er riktige.
b) Uendelig mange. $f'(x) = k(x + 1)(x - 2)$ med $k > 0$ gir toppunkt i $-1$ og bunnpunkt i $2$, og hver $k > 0$ og hver konstant i $f$ gir et nytt polynom. Med $k < 0$ blir det omvendt.
\fo{Sjekk 1} a) $f'(x) = 2x - 8$, bunnpunkt $(4, -13)$ \ b) Toppunkt \ c) Tangenten har stigningstall $f'(x)$. Er tallet positivt, stiger tangenten, og grafen følger tangenten like ved punktet, så $f$ stiger.

\fk{Del 2 — Største og minste verdi, optimering}
\fo{2.1} a) $f'(x) = 2x - 4$, $x = 2$ \ b) $f(0) = 0$, $f(2) = -4$, $f(5) = 5$ \ c) Største verdi $5$, minste verdi $-4$
\fo{2.2} a) $f'(x) = 3(x - 2)(x + 2)$. Verdier: $f(-3) = 9$, $f(-2) = 16$, $f(2) = -16$, $f(3) = -9$. Største verdi $16$, minste verdi $-16$.
b) $g'(x) = -2x + 4 = 0$ gir $x = 2$, utenfor $[3, 6]$. Verdier: $g(3) = 3$, $g(6) = -12$. Største verdi $3$, minste verdi $-12$.
\fo{2.3} a) Steg 1: $0 \le x \le 20$ \ b) Steg 2: $A(x) = 40x - 2x^2$ \ c) Steg 3: $A'(x) = 40 - 4x = 0$ gir $x = 10$ \ d) Steg 4: $A(0) = 0$, $A(10) = 200$, $A(20) = 0$.
Svar: Innhegningen får størst areal, $200$ m$^2$, når sidene vinkelrett på veggen er $10$ m og siden langs veggen er $20$ m.
\fo{2.4} a) Steg 1 og 2: Overflaten er $x^2 + 4xh = 300$, så $h = \frac{300 - x^2}{4x}$. Da er $V = x^2h = 75x - \frac{x^3}{4}$. Vi må ha $x > 0$ og $h > 0$, så $0 < x < 10\sqrt{3}$.
b) Steg 3: $V'(x) = 75 - \frac{3}{4}x^2 = 0$ gir $x^2 = 100$, så $x = 10$ ($x = -10$ er utenfor).
c) Steg 4: $V' > 0$ på $\langle 0, 10\rangle$ og $V' < 0$ på $\langle 10, 10\sqrt{3}\rangle$. $V$ vokser helt fram til $x = 10$ og avtar etterpå, så $V(10) = 500$ er største verdi. Da er $h = 5$.
Svar: Esken får størst volum, $500$ cm$^3$, når bunnen har side $10$ cm og høyden er $5$ cm.
\fo{2.5} Feilen er siste steg: «største verdi 5, minste verdi 4». Endepunktene er glemt. $f(0) = 0$ og $f(4) = 128 - 144 + 48 = 32$.
Riktig: største verdi $32$ (i $x = 4$) og minste verdi $0$ (i $x = 0$). Tallene $5$ og $4$ er bare lokalt topp og lokalt bunn.
\fo{2.6} a) $K'(t) = -1,5t^2 + 9t = 1,5t(6 - t)$, så $t = 0$ eller $t = 6$. $K(0) = 0$, $K(6) = 54$, $K(8) = 32$.
Størst: $54$ mg/L etter $6$ timer.
b) $K(1) = 4$, $K(6) = 54$, $K(8) = 32$. Laveste konsentrasjon er $4$ mg/L, i endepunktet $t = 1$. Der er $K'(1) = 7,5 \ne 0$: $K$ stiger videre, så grafen er ikke vannrett i endepunktet.
c) $K'(t) < 0$ for $t > 6$, så konsentrasjonen avtar fra $t = 6$ til $t = 8$ timer.
\fo{2.7} a) $f(x) - 2 = (x + 1)^2(x - 2) \le 0$ for $x \le 2$, så $f(x) \le 2$ for $x \le 2$, med likhet i $x = -1$ og $x = 2$. For $b < -1$ stiger $f$ på $[-2, b]$ og største verdi er $f(b) < 2$. For $-1 \le b \le 2$ er største verdi $2$. For $b > 2$ er $f(b) > 2$. Svar: $-1 \le b \le 2$.
b) Alltid. Er $f$ voksende, er $f(x) \le f(b)$ for alle $x$ i intervallet. Største verdi er i høyre endepunkt.
\fo{2.8} a) $f'(x) = 3x^2 - 27$ inneholder ikke $k$. Tallet $k$ flytter hele grafen opp eller ned, så alle $y$-verdier endres likt og rekkefølgen er den samme.
b) $f'(x) = 3(x - 3)(x + 3)$, kandidat $x = 3$. $f(0) = k$, $f(3) = -54 + k$, $f(6) = 54 + k$. Største verdi er $54 + k = 100$, så $k = 46$. Minste verdi er $-54 + 46 = -8$.
\fo{Sjekk 2} a) Største verdi $-5$ (i $x = 1$ og $x = 5$), minste verdi $-9$ (i $x = 3$) \ b) $A(x) = x(12 - x)$, $0 \le x \le 12$
c) Endepunktet kan ha den største eller minste verdien uten at grafen er vannrett der. Derfor må vi sammenligne verdiene.

\fk{Ekstra utfordring}
\fo{U1} a) $f(1) = 0$. Sett av en tom kolonne (tenk $0x^2$). Skriv produktet med motsatt fortegn og legg sammen, det er det samme som å trekke fra:
{\small\polylongdiv{x^3+0x^2-3x+2}{x-1}}
Dermed er $x^3 - 3x + 2 = (x - 1)(x^2 + x - 2) = (x - 1)^2(x + 2)$.
b) $f'(x) = 3(x - 1)(x + 1)$, fortegn $+, -, +$. Toppunkt $(-1, 4)$ og bunnpunkt $(1, 0)$. Nullpunkter $x = -2$ og $x = 1$. Skjæring med $y$-aksen: $(0, 2)$.
c) I dobbelt nullpunkt berører grafen $x$-aksen og snur. Her er bunnpunktet $(1, 0)$ også nullpunktet: $f(1) = 0$ og $f'(1) = 0$.
\fo{U2} a) Steg 1: $0 \le x \le 1500$. Steg 2: $F(x) = -0,02x^2 + 40x - 8000$. Steg 3: $F'(x) = -0,04x + 40 = 0$ gir $x = 1000$. Steg 4: $F(0) = -8000$, $F(1000) = 12\,000$, $F(1500) = 7000$.
Svar: Fortjenesten er størst, $12\,000$ kr, når bedriften selger $1000$ enheter.
b) Flere svar er riktige. $F(0) = -8000$ er faste kostnader, og $F(x) < 0$ for $x$ over ca. $1775$. Eksempel: Prisen må kanskje senkes ved høyt salg.
\fo{U3} Steg 1: $0 \le x \le 2\sqrt{3}$. Steg 2: $A(x) = 2x(12 - x^2) = 24x - 2x^3$. Steg 3: $A'(x) = 24 - 6x^2 = 0$ gir $x = 2$. Steg 4: $A(0) = 0$, $A(2) = 32$, $A(2\sqrt{3}) = 0$.
Svar: Rektangelet får størst areal, $32$, når bredden er $4$ og høyden er $8$.
\fo{U4} a) Utskrift: \texttt{3.92 1056.3052}. Verdien $x = \texttt{i}/100$ går fra $0$ til $10$ i steg på $0,01$, som er hele intervallet $[0, 10]$.
b) Eksakt: $x = \frac{25 - 5\sqrt{7}}{3} \approx 3,923748$ og $V \approx 1056,305895$. Feilen i $x$ er ca. $0,0037$ (steget er $0,01$). Feilen i $V$ er ca. $0,0007$, mye mindre. Nær toppunktet er $V' \approx 0$, så en liten feil i $x$ gir en mye mindre feil i $V$.
c) Endringer: \texttt{range(0, 10001)}, \texttt{x = i / 1000} og \texttt{round(beste\_x, 3)}. Utskrift: \texttt{3.924 1056.3059}.
\fo{U5} a) $f'(x) = 3x^2 + a$. Er $a > 0$, er $f'(x) > 0$ for alle $x$, så ingen ekstremalpunkter. Er $a = 0$, er $f'(x) = 3x^2 \ge 0$ uten fortegnsskifte. Er $a < 0$, er $f'(x) = 0$ for $x = \pm\sqrt{-a/3}$ og fortegnet skifter, så det er to ekstremalpunkter. Svar: $a \ge 0$.
b) $a = 0$. Terrassepunktet er $(0, 0)$.
\fo{U6} $f'(x) = 3px^2 + 2qx + r$ er et andregradspolynom med diskriminant $D = 4q^2 - 12pr$. Er $D > 0$, har $f'$ to nullpunkter, og parabelen krysser $x$-aksen i begge. Fortegnet skifter to ganger: ett toppunkt og ett bunnpunkt, altså $2$ ekstremalpunkter.
Er $D < 0$, har $f'$ ingen nullpunkter og samme fortegn overalt: $0$ ekstremalpunkter. Er $D = 0$, berører parabelen $x$-aksen uten å skifte fortegn: terrassepunkt, $0$ ekstremalpunkter. Altså $0$ eller $2$.
\fo{U7} $f'(x) = k(x + 1)(x - 1) = kx^2 - k$, så $f(x) = \frac{k}{3}x^3 - kx + c$. Toppunktet gir $f(-1) = \frac{2k}{3} + c = 7$ og bunnpunktet gir $f(1) = -\frac{2k}{3} + c = 3$. Differansen gir $\frac{4k}{3} = 4$, så $k = 3$, og da $c = 5$. Svar: $f(x) = x^3 - 3x + 5$. Fortegn for $f'$: $+, -, +$. Det finnes ikke flere: $k$ og $c$ er bestemt av de to punktene.
